% Nouveau EGA V (cote 25): archive cote 25, site page 116 (PDF page 117).
% typescript. Draft compared: none.
% AI-transcribed by Claude Opus 5.5 (claude-opus-5-5, adaptive thinking, effort high), Message Batches API, batch ega25b (msgbatch_01A7n5YuXubUPNK1zsepi8RG), 2026-11-13; not yet checked by eye.
\scanpage{25/116}
\mspage{19}

points fixes, elle est déterminée \xout{xxixxxxxxx} par le morphisme f:X $\longrightarrow$ P
associé. Noter d'ailleurs qu'on \xout{xx}obtient évidemment ainsi une correspon
dance biunivoque entre systêmes linéaires de diviseurs de X/S paramétré
par $\hand{\check{P}}$ , et morphismes f:X $\longrightarrow$ P tel que \xout{xx}$(fx_S\mathrm{id})^{-1}(\underline{H})$ soit un
diviseur\deleted{ de Cartier} \xout{xxx} relatif sur $Xx_S\check{P}$ par rapport à $\check{P}$ . Ce\over{\unsure{ci}}{tte} \deleted{se} \inserted{condition}
\inserted{se} vérifie aussi fibre par fibre, et on obtient :

\uline{Proposition} \over{\unsure{1}}{3}.1.On a une correspondance biunivoque entre les \xout{morphismes}
systêmes linéaires \inserted{sans points fixes} de diviseurs sur X/S paramétrés par $\hand{\check{P}}$ , et les morphis
mes f: X $\longrightarrow$ P ayant la propriété suivante: pour tout $s\in S$, d\over{e}{é}signant
par $\bar{k}$ une clôture algébrique de k(s), et pour tout cycle premier associé
X' de $X_{\bar{k}}$ , $f(X')\subset P_{\bar{k}}$ n'est contenu dans aucun hyperplan de $P_{\bar{k}}$ .
\xout{(xx}\hand{(}Lorsque X est à fibres géométriquement \xout{xxxxxxx} \inserted[typed]{intègres}, ceci \xout{xexsimplifie}
\xout{xx} s'énoncé simplement en disant que $f(X_{\bar{k}})$ n'est contenu dans aucun
hyperplan de $P_{\bar{k}}$ \hand{)}.

En général (i.e. \xout{xx} lorsque Z $\neq \emptyset$ ) on ne peut plus affirmer que
la connaissance de f détermine \xout{D}la famille de diviseurs. Le cas le plus
trivial est celui où $\hand{\check{P}}$ = S \inserted{est de dim relative 0}, alors la donnée d'un systême linéaire de\hand{.}
diviseurs de X/S paramétrée par S équivaut à celle d'un diviseur de
Cartier relatif D sur X par rapport à S; \xout{Txxxxlxxppx} le morphisme
associé est la projection X-D $\longrightarrow$ S , et on voit \xout{(mêmexxxx} que la connais
sance de ce morphisme (y compris celle de son domaine de définition)
ne détermine pas D. Dans ce cas également, supposant pour simplifier X
réduit, le domaine de définition de \xout{xxxmorphisme} f \xout{xlxsxxpas} considéré
comme application rationelle de X dans P = S n'est pas X-D , mais X.
Pour éliminer ce genre de phénomènes déplaisants, il y a lieu de se bor-%
ner aux systêmes \xout{h}linéaires de diviseurs \uline{``sans composante fixe''} . De
façon générale, si S=Spec(k), \unsure{étaht|étant} donné une famille (non néc linéaire)

\margin{116}
